
\documentclass[handout]{beamer}

\usetheme{Warsaw}

\usepackage[brazil]{babel}
\usepackage[latin1]{inputenc}
\usepackage{times}
\usepackage[T1]{fontenc}
\usepackage{epsfig}
\usepackage{listings}
\usepackage{amsfonts}
\usepackage{multirow}
\usepackage{multicol}
\usepackage{tikz}
\usepackage{textpos}
\usepackage{framed, multicol}
\usepackage{wrapfig}
\usepackage{array}
\usepackage{booktabs}
\usepackage{verbatim}
\usepackage{tikz}
\usepackage{tikz-uml}


%\usepackage{dcolumn}
%\newcolumntype{.}{D{.}{.}{-1}}
%\newcolumntype{d}[1]{D{.}{.}{#1}}
%\usetheme{Berkeley}

\newcommand{\be}{\begin{enumerate}[<+->]}
\newcommand{\ee}{\end{enumerate}}

\newcommand{\bq}{\begin{quote}}
\newcommand{\eq}{\end{quote}}

\newcommand{\bd}{\begin{description}[<+->]}
\newcommand{\ed}{\end{description}}

\newcommand{\bi}{\begin{itemize}[<+->]}
\newcommand{\ei}{\end{itemize}}

\newcommand{\bbl}{\begin{block}{Exemplo: Tipos primitivos do Pascal}
        \scriptsize
        \centering
        }
\newcommand{\ebl}{\end{block}}




%%%%%%%%%%%%%%%%%%%%%%%%%%%%%%%%%%%%%%%%%%%%%%%%%%%%%%%%%%%%%%%%
% Configuração do Tiks
%%%%%%%%%%%%%%%%%%%%%%%%%%%%%%%%%%%%%%%%%%%%%%%%%%%%%%%%%%%%%%%%
\usetikzlibrary{arrows}

\tikzstyle{bloco}=[draw, minimum size=2em] %fill=blue!20,
\tikzstyle{linha} = [pin edge={to-,thin,black}]



\title[Programa��o Modular]
{%
    Padr�es de projeto comportamentais%
}
\author[Prof. Hugo de Paula]
{
    Prof.~Hugo~de~Paula
}
\institute[DCC / PUC Minas]
{\includegraphics[width=4cm]{puc_engesoft_logo.png} \\
    \textsc{Pontif�cia Universidade Cat�lica de Minas Gerais}\\
    Curso de Engenharia de Software
}

\date[]{}

\lstset{language=Java,
   basicstyle=\scriptsize,
   commentstyle=\color{red},
   showstringspaces=false,
   numbers=none,
   numberstyle=\tiny}

\begin{document}


\selectlanguage{brazil}

\begin{frame}
   \titlepage
\end{frame}

%\addtobeamertemplate{frametitle}{}{%
%    \begin{tikzpicture}[remember picture,overlay]
%    \node[anchor=north east,yshift=2pt] at (current page.north east) {\epsfig{file=puclogo_small_bw,width=1.2cm}};
%    \end{tikzpicture}}


%\addtobeamertemplate{frametitle}{}{%
    %\begin{tikzpicture}[node distance=0cm, remember picture, overlay, every node/.style={inner sep=0,outer sep=0, node distance=0cm, baseline=0cm}]
    %\node[anchor=north east] at (current page.north east) {\epsfig{file=puclogo_small_bw,width=1cm}};
    %\end{tikzpicture}}


%\logo{\includegraphics[height=0.8cm]{puclogo_small_bw.pdf}\vspace{220pt}}


\begin{frame}
   \frametitle{Sum�rio}
   \tableofcontents[pausesections]
\end{frame}

%\AtBeginSection[] % Do nothing for \section*
%{
%\begin{frame}<beamer>
%\frametitle{Outline}
%\tableofcontents[currentsection]
%\end{frame}}

% duas   linhas 1.48cm
% tres   linhas 1.75cm
% quatro linhas 2.02cm
% cinco  linhas 2.26cm
\addtobeamertemplate{frametitle}{}{%
    \begin{textblock*}{10mm}(.8785\textwidth,-1.75cm)
        \includegraphics[height=0.97cm]{puc_engesoft_logo.png}%\includegraphics[height=1cm]{puclogo_small_bw.pdf}
    \end{textblock*}
}
%\addtobeamertemplate{frametitle}{}{%
%   \begin{textblock*}{10mm}(.9945\textwidth,-2.02cm)
%    \includegraphics[height=1cm]{puclogo_small_bw.pdf}
%   \end{textblock*}
%}

\section{Padr�es Comportamentais}

\begin{frame}{Padr�es de projeto comportamentais}

\begin{block}{Padr�es Comportamentais}
Organizam a comunica��o entre objetos de forma clara e eficiente.
\end{block}

\begin{itemize}
\item Exemplos mais comuns: Chain of Responsibility, Observer (Model-View-Controller), iterator, visitor, template e strategy.
\item Exemplos menos comuns: Interpreter e Mediator.


\end{itemize}

\end{frame}

\section{Chain of Responsibility}

\begin{frame}{Chain of Responsibility}

\begin{block}{Chain of Responsibility}
Descreve uma forma de organizar os objetos em cadeia (como na forma de uma corrente).
\end{block}

\begin{itemize}
\item �til quando alguns comandos devem ser tratados por diferentes objetos, cada um deles passando o comando ao pr�ximo objeto na forma de uma corrente.

\item Cada receptor de uma requisi��o cont�m uma refer�ncia a outro receptor.
\end{itemize}
\end{frame}


\subsection{Chain of Responsibility: exemplo Logger}

\begin{frame}{Chain of Responsibility: exemplo Logger}

\begin{center}
\includegraphics[width=.75\textwidth,angle=270]{chainLogger.pdf}
\end{center}

\end{frame}


\begin{frame}[fragile]{Chain of Responsibility: exemplo Logger}

\begin{lstlisting}
public abstract class Logger {
  public static int INFO = 1;
  public static int DEBUG = 2;
  public static int ERRO = 3;
  protected int nivel;
  // proximo elemento na cadeia de responsabilidades
  protected Logger next;

  public void setNext(Logger next) {
    this.next = next;
  }

  public void log(int nivel, String msg) {
    if (this.nivel <= nivel) {
      escrever(mensagem);
    }
    if (proximoLogger != null) {
      proximoLogger.log(nivel, msg);
    }
  }
  abstract protected void write(String msg);
}
\end{lstlisting}
\end{frame}


\begin{frame}[fragile]{Chain of Responsibility: exemplo Logger}

\begin{lstlisting}[basicstyle=\tiny]
public class LoggerConsole extends Logger {
  public LoggerConsole(int nivel) {
    this.nivel = nivel;
  }
  @Override
  protected void write(String msg) {
    System.out.println("Console Padrao::Logger: " + msg);
  }
}

public class LoggerErro extends Logger {
  public LoggerErro(int nivel) {
    this.nivel = nivel;
  }
  @Override
  protected void write(String msg) {
    System.out.println("Console de Error::Logger: " + msg);
  }
}

public class LoggerArquivo extends Logger {
  public LoggerArquivo(int nivel) {
    this.nivel = nivel;
  }
  @Override
  protected void write(String msg) {
    System.out.println("Arquivo::Logger: " + msg);
  }
}
\end{lstlisting}
\end{frame}



\begin{frame}[fragile]{Chain of Responsibility: exemplo Logger}

\begin{lstlisting}

public class Cliente {

  private static Logger getCadeiaDeLoggers() {
    Logger errorLogger = new LoggerErro(Logger.ERRO);
    Logger fileLogger = new LoggerArquivo(Logger.DEBUG);
    Logger consoleLogger = new LoggerConsole(Logger.INFO);
    errorLogger.setNext(fileLogger);
    fileLogger.setNext(consoleLogger);
    return errorLogger;
  }

  public static void main(String[] args) {
    Logger loggerChain = getCadeiaDeLoggers();
    loggerChain.log(Logger.INFO, "Esta � uma informa��o de evento.");
    loggerChain.log(Logger.DEBUG, "Esta � uma informa��o de Debug.");
    loggerChain.log(Logger.ERRO, "Esta � uma mensagem de erro.");
  }
}
\end{lstlisting}
\end{frame}

\section{Observer}

\begin{frame}{Observer}

    \begin{block}{Observer}
        Rela��o de um-para-muitos entre objetos de tal forma que quando um objeto � modificado os objetos dependentes s�o notificados automaticamente.
    \end{block}
\end{frame}


\subsection{Solu��o geral}

\begin{frame}{Observer: solu��o geral}
    \begin{center}\tiny
        \begin{tikzpicture}
        \umlclass[type=abstract, font=\fontsize{8}{8}\selectfont]{Observer}{-- subject : Subject}{+ \textit{notify()}}
        \umlclass[font=\fontsize{8}{8}\selectfont,y=-2]{ConcreteObserver}{}{+ notify()}
        \umlclass[font=\fontsize{8}{8}\selectfont,x=5,y=-0.7]{Subject}{-- observers : List<Observer> \\ -- state : State }{ + getState() : State \\ + setState(: State) : void \\ + attach(: Observer) : void \\ + notifyAllObservers() : void}
        \umlVHVinherit{ConcreteObserver}{Observer}
        \umluniassoc[font=\fontsize{8}{8}\selectfont,arg2=notify,pos2=0.5]{Subject}{Observer}
        \end{tikzpicture} \\[5mm]
    \end{center}

    \begin{itemize}
        \item \textbf{Subject}: permite a conex�o de observadores e possui o estado que ser� observado.
        \item \textbf{Observer}: se conecta a um Subject e declara o m�todo que � executado no momento da notifica��o.
        \item \textbf{ConcreteObserver}: implementa o m�todo de notifica��o.
    \end{itemize}

\end{frame}



\subsection{Exemplo: Observer}

\begin{frame}{Exemplo: Observer}
    \begin{center}\tiny
        \begin{tikzpicture}
        \umlclass[type=abstract, font=\fontsize{8}{8}\selectfont,y=1.5]{Observer}{-- subject : Subject}{+ \textit{update()}}
        \umlclass[font=\fontsize{8}{8}\selectfont,x=-3, y=-2]{BinaryObserver}{}{+ update()}
        \umlclass[font=\fontsize{8}{8}\selectfont,x=0, y=-2]{OctalObserver}{}{+ update()}
        \umlclass[font=\fontsize{8}{8}\selectfont,x=3, y=-2]{HexaObserver}{}{+ update()}
        \umlVHVinherit{BinaryObserver}{Observer}
        \umlVHVinherit{OctalObserver}{Observer}
        \umlVHVinherit{HexaObserver}{Observer}
        \umlclass[font=\fontsize{8}{8}\selectfont,x=5, y=1.5]{Subject}{-- observers : List<Observer> \\ -- state : int }{ + getState() : int \\ + setState(: State) : void \\ + attach(: Observer) : void \\ + notifyAllObservers() : void}
        \umluniassoc[font=\fontsize{8}{8}\selectfont,arg2=notify, pos2=0.5]{Subject}{Observer}
        \end{tikzpicture} \\[5mm]
    \end{center}
\end{frame}




\end{document}










